% GENERATED FILE. Edit canonical lessons, then run npm run generate:books.
% canonical-source-hash: fnv1a64:1cb2654ade809ddb
% canonical-lessons: BN-C386-kuli, BN-C386-sutkes, BN-C386-troli, BN-C386-namano, BN-C386-bohon-kora

\chapter{Porters and Luggage}
\label{ch:bn-porters-and-luggage}
\hlchaptermodality{386}

\begin{chapteropening}
\textbf{By the end of this chapter:} \emph{I can get help to carry and unload my bags.}

Complete the last lesson of chapter 386: I can get help to carry and unload my bags.
\end{chapteropening}

\section[\texorpdfstring{\bn{কুলি} (kuli)}{কুলি (kuli)}]{\bn{কুলি} (kuli) — a porter}

\label{lesson:BN-C386-kuli}



\begin{quote}
Before the new one: say the Bengali for to wait, then the Bengali for a bench.
\end{quote}

\subsection*{You'll want to know: \bn{কুলি}}
\textbf{\bn{কুলি}} — \emph{kuli} — \textquotedblleft{}a porter\textquotedblright{}.

\begin{quote}
\bn{কুলি} \bn{আমার} \bn{ভারী} \bn{মালপত্র} \bn{তুলল।}
\end{quote}

\begin{quote}
\emph{kuli āmār bhārī mālpôtrô tullô.}
\end{quote}

\begin{quote}
The porter lifted my heavy luggage.
\end{quote}

At a big station the porters wear red shirts and carry loads on the head.

\begin{cousinweb}
\textbf{\bn{কুলি}} (\emph{kuli}) has an origin that is not settled: it is often traced to the South Indian word for wages, Tamil \emph{kūli} (DEDR 1905), so that a porter is someone paid by the job.
\end{cousinweb}

\begin{grammarlens}[title={What you've built}]
One more word for getting around.
\end{grammarlens}

\subsection*{Your turn}
\begin{itemize}
\raggedright
  \item \emph{Say it:} \emph{kuli}
  \item \emph{Say it:} \emph{kuli}, once more
  \item \emph{Say it:} \emph{kuli āmār bhārī mālpôtrô tullô.} — The porter lifted my heavy luggage.
  \item \emph{Recall:} say the Bengali for to wait, then the Bengali for a bench, then say \emph{kuli} again
\end{itemize}

\begin{tcolorbox}[breakable,colback=teal!4,colframe=teal!35!black,title={Before you move on}]
What does \bn{কুলি} mean? (\textquotedblleft{}A porter\textquotedblright{}.) Say it once more.
\end{tcolorbox}
\section[\texorpdfstring{\bn{সুটকেস} (suṭkes)}{সুটকেস (suṭkes)}]{\bn{সুটকেস} (suṭkes) — a suitcase}

\label{lesson:BN-C386-sutkes}



\begin{quote}
Before the new one: say the Bengali for a bench, then the Bengali for a porter.
\end{quote}

\subsection*{You'll want to know: \bn{সুটকেস}}
\textbf{\bn{সুটকেস}} — \emph{suṭkes} — \textquotedblleft{}a suitcase\textquotedblright{}.

\begin{quote}
\bn{আমি} \bn{সুটকেসটা} \bn{খাটের} \bn{নিচে} \bn{রাখলাম।}
\end{quote}

\begin{quote}
\emph{āmi suṭkesṭā khāṭer niche rākhlām.}
\end{quote}

\begin{quote}
I put the suitcase under the bed.
\end{quote}

\begin{cousinweb}
\textbf{\bn{সুটকেস}} (\emph{suṭkes}) is borrowed from English \emph{suitcase}, a case for a suit of clothes, with the English t said as the hard Bengali \bn{ট}.
\end{cousinweb}

\begin{grammarlens}[title={What you've built}]
One more word for getting around.
\end{grammarlens}

\subsection*{Your turn}
\begin{itemize}
\raggedright
  \item \emph{Say it:} \emph{suṭkes}
  \item \emph{Say it:} \emph{suṭkes}, once more
  \item \emph{Say it:} \emph{āmi suṭkesṭā khāṭer niche rākhlām.} — I put the suitcase under the bed.
  \item \emph{Recall:} say the Bengali for a bench, then the Bengali for a porter, then say \emph{suṭkes} again
\end{itemize}

\begin{tcolorbox}[breakable,colback=teal!4,colframe=teal!35!black,title={Before you move on}]
What does \bn{সুটকেস} mean? (\textquotedblleft{}A suitcase\textquotedblright{}.) Say it once more.
\end{tcolorbox}
\section[\texorpdfstring{\bn{ট্রলি} (ṭrôli)}{ট্রলি (ṭrôli)}]{\bn{ট্রলি} (ṭrôli) — a trolley}

\label{lesson:BN-C386-troli}



\begin{quote}
Before the new one: say the Bengali for a porter, then the Bengali for a suitcase.
\end{quote}

\subsection*{You'll want to know: \bn{ট্রলি}}
\textbf{\bn{ট্রলি}} — \emph{ṭrôli} — \textquotedblleft{}a trolley\textquotedblright{}.

\begin{quote}
\bn{বিমানবন্দরে} \bn{আমরা} \bn{একটা} \bn{ট্রলি} \bn{নিলাম।}
\end{quote}

\begin{quote}
\emph{bimānbôndôre āmrā ekṭā ṭrôli nilām.}
\end{quote}

\begin{quote}
We took a trolley at the airport.
\end{quote}

\begin{cousinweb}
\textbf{\bn{ট্রলি}} (\emph{ṭrôli}) is borrowed from English \emph{trolley}, the wheeled cart of airports and stations.
\end{cousinweb}

\begin{grammarlens}[title={What you've built}]
One more word for getting around.
\end{grammarlens}

\subsection*{Your turn}
\begin{itemize}
\raggedright
  \item \emph{Say it:} \emph{ṭrôli}
  \item \emph{Say it:} \emph{ṭrôli}, once more
  \item \emph{Say it:} \emph{bimānbôndôre āmrā ekṭā ṭrôli nilām.} — We took a trolley at the airport.
  \item \emph{Recall:} say the Bengali for a porter, then the Bengali for a suitcase, then say \emph{ṭrôli} again
\end{itemize}

\begin{tcolorbox}[breakable,colback=teal!4,colframe=teal!35!black,title={Before you move on}]
What does \bn{ট্রলি} mean? (\textquotedblleft{}A trolley\textquotedblright{}.) Say it once more.
\end{tcolorbox}
\section[\texorpdfstring{\bn{নামানো} (nāmāno)}{নামানো (nāmāno)}]{\bn{নামানো} (nāmāno) — to take down; to unload}

\label{lesson:BN-C386-namano}



\begin{quote}
Before the new one: say the Bengali for a suitcase, then the Bengali for a trolley.
\end{quote}

\subsection*{You'll want to know: \bn{নামানো}}
\textbf{\bn{নামানো}} — \emph{nāmāno} — \textquotedblleft{}to take down; to unload\textquotedblright{}.

\begin{quote}
\bn{অরুণ} \bn{তাক} \bn{থেকে} \bn{বাক্সটা} \bn{নামাল।}
\end{quote}

\begin{quote}
\emph{ôrun tāk theke bāksôṭā nāmālô.}
\end{quote}

\begin{quote}
Arun took the box down from the shelf.
\end{quote}

Past (he/she …): \bn{তিনি} \bn{বাক্সটা} \bn{নামালেন}, \emph{tini bāksôṭā nāmālen}.

The causing partner of \bn{নামা}, to get down: you \bn{নামো} from a bus, and you \bn{নামাও} a box from a shelf.

\begin{cousinweb}
\textbf{\bn{নামানো}} (\emph{nāmāno}) is the causing form of \textbf{\bn{নামা}}, \textquotedblleft{}to get down\textquotedblright{}, which CDIAL lists as Bengali \emph{nāmā}, \textquotedblleft{}to descend\textquotedblright{}, under Sanskrit \emph{namrá}, \textquotedblleft{}bent, bowed\textquotedblright{} (CDIAL 6964).
\end{cousinweb}

\begin{grammarlens}[title={What you've built}]
One more word for getting around.
\end{grammarlens}

\subsection*{Your turn}
\begin{itemize}
\raggedright
  \item \emph{Say it:} \emph{nāmāno}
  \item \emph{Say it:} \emph{nāmāno}, once more
  \item \emph{Say it:} \emph{ôrun tāk theke bāksôṭā nāmālô.} — Arun took the box down from the shelf.
  \item \emph{Recall:} say the Bengali for a suitcase, then the Bengali for a trolley, then say \emph{nāmāno} again
\end{itemize}

\begin{tcolorbox}[breakable,colback=teal!4,colframe=teal!35!black,title={Before you move on}]
What does \bn{নামানো} mean? (\textquotedblleft{}To take down; to unload\textquotedblright{}.) Say it once more.
\end{tcolorbox}
\section[\texorpdfstring{\bn{বহন} \bn{করা} (bôhôn kôrā)}{বহন করা (bôhôn kôrā)}]{\bn{বহন} \bn{করা} (bôhôn kôrā) — to carry (a load); to bear}

\label{lesson:BN-C386-bohon-kora}



\begin{quote}
Before the new one: say the Bengali for a trolley, then the Bengali for to unload.
\end{quote}

\subsection*{You'll want to know: \bn{বহন} \bn{করা}}
\textbf{\bn{বহন} \bn{করা}} — \emph{bôhôn kôrā} — \textquotedblleft{}to carry (a load); to bear\textquotedblright{}.

\begin{quote}
\bn{ট্রলিটা} \bn{আমাদের} \bn{সব} \bn{মালপত্র} \bn{বহন} \bn{করল।}
\end{quote}

\begin{quote}
\emph{ṭrôliṭā āmāder sôb mālpôtrô bôhôn kôrlô.}
\end{quote}

\begin{quote}
The trolley carried all our luggage.
\end{quote}

Past (he/she …): \bn{সেই} \bn{ট্রাক} \bn{মালপত্র} \bn{বহন} \bn{করল}, \emph{sei ṭrāk mālpôtrô bôhôn kôrlô}.

A careful word, common in writing and on signs; in talk people say \bn{নিয়ে} \bn{যাওয়া}, to take along.

\begin{cousinweb}
\textbf{\bn{বহন} \bn{করা}} (\emph{bôhôn kôrā}) joins Sanskrit \emph{vahana}, \textquotedblleft{}bearing, carrying, conveying\textquotedblright{} (Monier-Williams), to \bn{করা}, to do.
\end{cousinweb}

\begin{grammarlens}[title={What you've built}]
One more word for getting around.
\end{grammarlens}

\subsection*{Your turn}
\begin{itemize}
\raggedright
  \item \emph{Say it:} \emph{bôhôn kôrā}
  \item \emph{Say it:} \emph{bôhôn kôrā}, once more
  \item \emph{Say it:} \emph{ṭrôliṭā āmāder sôb mālpôtrô bôhôn kôrlô.} — The trolley carried all our luggage.
  \item \emph{Recall:} say the Bengali for a trolley, then the Bengali for to unload, then say \emph{bôhôn kôrā} again
\end{itemize}

\begin{tcolorbox}[breakable,colback=teal!4,colframe=teal!35!black,title={Before you move on}]
What does \bn{বহন} \bn{করা} mean? (\textquotedblleft{}To carry (a load); to bear\textquotedblright{}.) Say it once more.
\end{tcolorbox}
